\documentclass[aspectratio=169,11pt]{beamer}

\usetheme{Madrid}
\usecolortheme{crane}
\setbeamertemplate{navigation symbols}{}

\definecolor{darkred}{RGB}{139,0,0}
\definecolor{brightred}{RGB}{220,20,60}
\setbeamercolor{title}{fg=white,bg=darkred}
\setbeamercolor{frametitle}{fg=white,bg=darkred}
\setbeamercolor{section in head/foot}{bg=brightred,fg=white}
\setbeamercolor{alerted text}{fg=darkred}

\usepackage[T1]{fontenc}
\usepackage[utf8]{inputenc}
\usepackage[french]{babel}
\usepackage{graphicx}
\usepackage{booktabs}
\usepackage{amsmath}
\usepackage{hyperref}

\title{Methodologie detaillee de reconstruction paleofeu}
\subtitle{CharAnalysis, lissage, agregations regionales et bootstrap}
\author{Projet Foret Boreale}
\institute{Synthese technique basee sur l'article principal et ses references methodologiques}
\date{\today}

\begin{document}

\begin{frame}
  \titlepage
\end{frame}

\begin{frame}{Plan}
  \tableofcontents
\end{frame}

\section{1. Objectif et donnees}

\begin{frame}{Question methodologique}
\textbf{Objectif:} reconstruire la dynamique du feu sur 12 000 ans a partir d'archives lacustres.

\vspace{0.3cm}
\textbf{Trois sorties cibles:}
\begin{itemize}
  \item \alert{RegBB}: biomasse brulee regionale.
  \item \alert{RegFRI}: intervalle de retour regional du feu.
  \item \alert{RegFS}: indice regional de taille/severite.
\end{itemize}

\textbf{Principe general:}
\begin{enumerate}
  \item Transformer le signal charbon en series temporelles comparables.
  \item Extraire les evenements (pics) et la frequence des feux.
  \item Agreger 6 lacs en un signal regional robuste.
  \item Quantifier l'incertitude par bootstrap (IC 90\%).
\end{enumerate}
\end{frame}

\begin{frame}{Jeu de donnees paleofeu (entree)}
\textbf{Sites:} 6 lacs (M14, M15, O6, O14, O15, O16).

\textbf{Materiel sedimentaire:}
\begin{itemize}
  \item carottes Livingstone (long terme),
  \item carottes Kajak-Brinkhurst (surface recente).
\end{itemize}

\textbf{Echelle de travail:}
\begin{itemize}
  \item sous-echantillon sedimentaire: 1 cm$^3$ par tranche de 1 cm,
  \item resolution temporelle finale: variable selon les modeles age-profondeur,
  \item fenetre d'analyse regionale: 0--12 000 cal. yr BP.
\end{itemize}
\end{frame}

\begin{frame}{Chronologie age-profondeur (fondement du calcul)}
\textbf{Recent:} datation \(^{210}\)Pb (modele CRS).

\textbf{Ancien:} datations \(^{14}\)C (AMS) sur macro-restes terrestres et/ou gyttja.

\textbf{Modelisation:}
\begin{itemize}
  \item calibration IntCal20,
  \item modele bayesien (type rbacon),
  \item interpolation des ages a 1 cm,
  \item CI autour des ages calibres.
\end{itemize}

\textbf{Pourquoi c'est critique:}
\begin{itemize}
  \item toute erreur de chronologie se propage dans CHAR,
  \item elle affecte directement detection des pics, FF, FRI et RegFS.
\end{itemize}
\end{frame}

\section{2. Extraction charbon et CHAR}

\begin{frame}{Extraction des macro-charbons}
\textbf{Pipeline laboratoire:}
\begin{enumerate}
  \item prelever 1 cm$^3$ de sediment,
  \item disperser avec solution NaPO\(_3\) a 3\%,
  \item tamiser a \alert{160 \(\mu\)m},
  \item imager les particules (grossissement 20x),
  \item mesurer nombre et aire des fragments.
\end{enumerate}

\textbf{Choix d'analyse:}
\begin{itemize}
  \item nombre et aire sont fortement correles,
  \item l'etude retient principalement \alert{l'aire du charbon}.
\end{itemize}
\end{frame}

\begin{frame}{Pourquoi le seuil 160 \(\mu\)m?}
\begin{itemize}
  \item Les fragments >160 \(\mu\)m sont associes a des sources de feu locales/subregionales.
  \item L'ordre de grandeur spatial est generalement \alert{1 a 30 km} autour du lac.
  \item Ce choix favorise la reconstruction des regimes de feu du paysage proche,
        plutot qu'un signal hemispherique diffus.
\end{itemize}

\vspace{0.2cm}
\textbf{Consequence:}
\begin{itemize}
  \item meilleure sensibilite aux evenements locaux,
  \item mais dependance plus forte a l'heterogeneite entre lacs.
\end{itemize}
\end{frame}

\begin{frame}{Fusion des carottes et serie continue par lac}
\textbf{Probleme:} deux types de carottes couvrent des portions temporelles differentes.

\textbf{Procedure:}
\begin{itemize}
  \item comparer motifs de nombre/aire dans les zones de recouvrement,
  \item aligner les segments,
  \item conserver prioritairement la partie recente KB (mieux contrainte en \(^{210}\)Pb),
  \item retirer les duplications Livingstone dans la zone de recouvrement.
\end{itemize}

\textbf{Sortie:} une serie charbon unique continue pour chaque lac.
\end{frame}

\begin{frame}{Conversion en CHAR}
\textbf{Entree brute:} concentration/aire charbon volumique
\[
C(t) \; [\mathrm{mm}^2\,\mathrm{cm}^{-3}]
\]

\textbf{Sortie standardisee:}
\[
CHAR(t) \; [\mathrm{mm}^2\,\mathrm{cm}^{-2}\,\mathrm{an}^{-1}]
\]

\textbf{Principe:}
\begin{itemize}
  \item conversion via la vitesse d'accumulation sedimentaire issue du modele age-profondeur,
  \item interpolation sur une grille temporelle harmonisee (resolution mediane).
\end{itemize}

\textbf{Interpretation:}
\begin{itemize}
  \item CHAR eleve = flux de charbon eleve,
  \item proxy de biomasse brulee locale plus importante.
\end{itemize}
\end{frame}

\begin{frame}{Pourquoi la vitesse de sedimentation compte?}
\textbf{Idee cle:} la conversion en CHAR depend directement de la chronologie et du taux d'accumulation sedimentaire.

\vspace{0.2cm}
\begin{itemize}
  \item Si les sediments s'accumulent vite, une meme quantite de charbon correspond a un flux annuel different.
  \item Si les sediments s'accumulent lentement, le meme volume couvre plus de temps.
  \item Le modele age-profondeur fournit la relation profondeur-temps qui permet cette conversion.
\end{itemize}

\vspace{0.2cm}
\textbf{Traduction simple:}
\begin{itemize}
  \item le charbon n'est pas seulement mesure par volume,
  \item il est converti en flux temporel grace au contexte sedimentaire.
\end{itemize}
\end{frame}

\begin{frame}{Exemple intuitif de conversion en CHAR}
\textbf{Cas 1: sedimentation rapide}
\begin{itemize}
  \item 1 cm de sediment se depose en 20 ans.
  \item Le charbon contenu dans cette tranche est rapporte a une duree courte.
\end{itemize}

\textbf{Cas 2: sedimentation lente}
\begin{itemize}
  \item 1 cm de sediment se depose en 100 ans.
  \item La meme tranche couvre plus de temps, donc le flux annuel calcule change.
\end{itemize}

\vspace{0.2cm}
\textbf{Conclusion:} deux echantillons ayant la meme aire de charbon peuvent donner des CHAR differents si leur vitesse de sedimentation differe.
\end{frame}

\section{3. CharAnalysis en detail}

\begin{frame}{Role de CharAnalysis}
\textbf{But:} separer, dans CHAR, le signal de fond (tendance lente) du signal d'evenements (pics).

\vspace{0.2cm}
\textbf{Sorties principales:}
\begin{itemize}
  \item dates d'evenements feu (pics),
  \item serie de biomasse brulee,
  \item base pour calcul de FF, FRI et RegFS.
\end{itemize}

\textbf{Logiciel mentionne:} CharAnalysis v1.1.
\end{frame}

\begin{frame}{Lecture intuitive: fond vs pics}
\textbf{Signal de fond (background):}
\begin{itemize}
  \item representer la dynamique lente du systeme,
  \item suivre les tendances multi-centennales,
  \item ne pas etre interprete comme une succession d'evenements individuels.
\end{itemize}

\textbf{Signal de pics (peak component):}
\begin{itemize}
  \item representer les deviations rapides au-dessus du fond,
  \item candidate a l'interpretation "evenement de feu".
\end{itemize}

\textbf{Message cle:} CharAnalysis ne detecte pas les feux sur CHAR brut, mais sur le residuel apres retrait du fond.
\end{frame}

\begin{frame}{Etape A -- Interpolation temporelle}
\textbf{Pourquoi:} les intervalles entre echantillons ne sont pas constants.

\textbf{Action:}
\begin{itemize}
  \item interpoler CHAR sur une grille temporelle reguliere,
  \item utiliser une resolution proche de la resolution mediane du site.
\end{itemize}

\textbf{Point de vigilance:}
\begin{itemize}
  \item une interpolation trop fine amplifie le bruit,
  \item trop grossiere masque les pics et sous-estime les evenements.
\end{itemize}
\end{frame}

\begin{frame}{Etape B -- Estimation du bruit de fond}
\textbf{Concept:}
\[
CHAR_{interp}(t) = CHAR_{background}(t) + CHAR_{peak}(t)
\]

\textbf{Background:}
\begin{itemize}
  \item composante basse frequence (variations lentes),
  \item obtenue par lissage fenetre/mobile robuste (selon parametrage CharAnalysis).
\end{itemize}

\textbf{Peak component:}
\[
CHAR_{peak}(t) = CHAR_{interp}(t) - CHAR_{background}(t)
\]
\end{frame}

\begin{frame}{Etape C -- Detection de pics}
\textbf{Principe:}
\begin{itemize}
  \item estimer la distribution du bruit residuel,
  \item definir un seuil statistique,
  \item detecter comme \alert{evenement} les maxima depassant ce seuil.
\end{itemize}

\textbf{Resultat:}
\begin{itemize}
  \item suite de dates de feu pour chaque lac,
  \item base de calcul de la frequence des feux (FF).
\end{itemize}

\textbf{Bonnes pratiques:}
\begin{itemize}
  \item tester la sensibilite aux parametres,
  \item comparer avec references dendro/archives historiques si disponibles,
  \item documenter explicitement seuils et fenetres choisis.
\end{itemize}
\end{frame}

\begin{frame}{Que veut dire \textit{pic significatif}?}
\textbf{Definition operationnelle:}
\begin{itemize}
  \item un pic est dit significatif s'il depasse un seuil estime a partir du bruit residuel,
  \item il est alors classe comme evenement plausible plutot que fluctuation aleatoire.
\end{itemize}

\textbf{Implication pratique:}
\begin{itemize}
  \item seuil trop bas \(\Rightarrow\) faux positifs (trop d'evenements),
  \item seuil trop haut \(\Rightarrow\) faux negatifs (evenements manques).
\end{itemize}

\textbf{Bonne pratique:} justifier le parametrage et faire une analyse de sensibilite.
\end{frame}

\begin{frame}{Comment choisir le seuil statistique?}
\textbf{Principe:} compromis entre faux positifs et faux negatifs.

\textbf{Procedure pratique (CharAnalysis):}
\begin{enumerate}
  \item Estimer la distribution du bruit residuel de \(CHAR_{peak}\),
  \item definir un seuil sur cette distribution (souvent percentile eleve),
  \item tester la sensibilite sur plusieurs seuils.
\end{enumerate}

\textbf{Regle de travail:}
\begin{itemize}
  \item commencer conservateur (ex. percentile 99\%),
  \item comparer avec 97.5\% et 95\%,
  \item retenir le seuil qui stabilise FF sans sur-detection.
\end{itemize}

\textbf{A reporter dans la methode:} seuil final, variantes testees, impact sur le nombre de pics et la FF regionale.
\end{frame}

\begin{frame}{Sortie CharAnalysis -> entree de FF}
\textbf{Sortie directe de CharAnalysis:}
\begin{itemize}
  \item liste datee des pics par lac,
  \item intensite relative des pics.
\end{itemize}

\textbf{Transition methodologique:}
\begin{itemize}
  \item ces dates de pics sont transformees en frequence de feu,
  \item via une estimation de densite a noyau (bande passante 500 ans).
\end{itemize}

\textbf{Transition orale (a dire):} \alert{"Les dates de pics permettent maintenant de calculer la frequence des feux."}
\end{frame}

\begin{frame}[fragile]{Pseudo-code conceptuel CharAnalysis}
\footnotesize
\begin{verbatim}
Input: charcoal_area_by_depth, age_depth_model
1) Convert depth -> age
2) Compute CHAR(age)
3) Interpolate CHAR on regular time grid
4) Smooth CHAR to estimate background component
5) Residual = CHAR_interp - CHAR_background
6) Fit noise model on residual component
7) Detect peaks above threshold
Output: fire-event dates, CHAR background, peak series
\end{verbatim}

\normalsize
\textbf{Important:} le detail exact des parametres (fenetres, seuil) doit etre reporte dans l'annexe methodologique du projet.
\end{frame}

\section{4. Frequence, FRI, taille/severite}

\begin{frame}{Frequence des feux (FF): entree et objectif}
\textbf{Entree:} dates de pics detectees par CharAnalysis pour chaque lac.

\textbf{Pourquoi ne pas juste compter par fenetre fixe?}
\begin{itemize}
  \item le comptage par classes peut etre instable aux bords de fenetre,
  \item la densite a noyau fournit une estimation continue de la frequence,
  \item elle exploite l'information temporelle de chaque evenement.
\end{itemize}

\textbf{Sortie:}
\[
FF(t)\; [\mathrm{feux}\,\mathrm{an}^{-1}]
\]
\end{frame}

\begin{frame}{Estimation de FF par densite a noyau (h = 500 ans)}
\textbf{Forme generale:}
\[
FF(t) = \frac{1}{h}\sum_i K\!\left(\frac{t-t_i}{h}\right),\qquad h=500\ \text{ans}
\]
\begin{itemize}
  \item \(t_i\): dates des pics,
  \item \(K\): noyau de lissage,
  \item \(h\): bande passante qui controle l'echelle de lissage.
\end{itemize}

\textbf{Effet de \(h=500\):}
\begin{itemize}
  \item attenue le bruit haute frequence,
  \item conserve les tendances centennales a millenaires,
  \item rend les courbes comparables entre lacs.
\end{itemize}
\end{frame}

\begin{frame}{Bande passante: modifiable ou non?}
\textbf{A quoi correspond \(h\)?}
\begin{itemize}
  \item \(h\) est la largeur de lissage temporel de la densite a noyau,
  \item petit \(h\): plus de detail mais plus de bruit,
  \item grand \(h\): courbe plus stable mais perte des variations rapides.
\end{itemize}

\textbf{Est-ce modifiable?} Oui, c'est un choix methodologique.

\textbf{Comment modifier?}
\begin{itemize}
  \item recalculer FF avec plusieurs valeurs (ex. 300, 500, 700 ans),
  \item comparer la stabilite des tendances regionales,
  \item retenir la valeur qui limite le bruit sans effacer le signal utile.
\end{itemize}

\textbf{Pourquoi ne pas modifier en cours d'analyse?}
\begin{itemize}
  \item pour garder la comparabilite entre periodes et entre sites,
  \item pour eviter un reglage opportuniste des resultats,
  \item pour conserver une methode reproductible et defendable.
\end{itemize}
\end{frame}

\begin{frame}{Comment lire FF?}
\textbf{Interpretation directe:}
\begin{itemize}
  \item FF eleve \(\Rightarrow\) feux plus frequents (intervalle moyen plus court),
  \item FF faible \(\Rightarrow\) feux plus espaces (intervalle moyen plus long).
\end{itemize}

\textbf{Exemple numerique:}
\begin{itemize}
  \item si \(FF=0.01\) feu/an, on attend en moyenne 1 feu tous les 100 ans,
  \item si \(FF=0.004\) feu/an, on attend en moyenne 1 feu tous les 250 ans.
\end{itemize}

\textbf{Transition:} une fois FF estime, on deduit l'intervalle de retour \(FRI=1/FF\).
\end{frame}

\begin{frame}{De FF a FRI}
\textbf{Definition:}
\[
FRI(t) = \frac{1}{FF(t)}
\]

\textbf{Unites:}
\begin{itemize}
  \item FF en feux/an,
  \item FRI en annees entre feux.
\end{itemize}

\textbf{Lecture:}
\begin{itemize}
  \item FF fort \(\Rightarrow\) FRI court (feux rapproches),
  \item FF faible \(\Rightarrow\) FRI long (feux espaces).
\end{itemize}

\textbf{Regional:}
\[
RegFRI(t) = \frac{1}{RegFF(t)}
\]
\end{frame}

\begin{frame}{Biomasse brulee regionale (RegBB)}
\textbf{Construction:}
\begin{itemize}
  \item convertir chaque serie lacustre en proxy BB,
  \item harmoniser l'echelle temporelle,
  \item agreger les 6 series en signal regional.
\end{itemize}

\textbf{Interpretation:}
\begin{itemize}
  \item RegBB eleve: combustion regionale soutenue,
  \item RegBB bas: combustion regionale reduite.
\end{itemize}
\end{frame}

\begin{frame}{Indice taille/severite (RegFS)}
\textbf{Definition dans l'etude:}
\[
RegFS(t) = \frac{RegBB(t)}{RegFF(t)}
\]

\textbf{Intuition:}
\begin{itemize}
  \item c'est une quantite moyenne de biomasse brulee par evenement,
  \item elle combine information de combustion totale et nombre d'evenements.
\end{itemize}

\textbf{Lecture:}
\begin{itemize}
  \item RegFS haut: feux moyens plus grands et/ou plus severes,
  \item RegFS bas: feux moyens plus petits et/ou moins severes.
\end{itemize}
\end{frame}

\section{5. Aggregation regionale et bootstrap}

\begin{frame}{Aggregation multi-sites}
\textbf{Pourquoi agreger 6 lacs?}
\begin{itemize}
  \item reduire l'effet des idiosyncrasies locales,
  \item extraire un signal de regime regional,
  \item tester la coherence inter-sites.
\end{itemize}

\textbf{Strategie:}
\begin{enumerate}
  \item calculer les metriques site par site,
  \item aligner sur la meme grille temporelle,
  \item combiner en moyenne regionale (RegBB, RegFF, RegFS),
  \item derivation de RegFRI.
\end{enumerate}
\end{frame}

\begin{frame}{Bootstrap: principe}
\textbf{Objectif:} propager l'incertitude inter-sites vers les indicateurs regionaux.

\textbf{Configuration de l'etude:}
\begin{itemize}
  \item resampling des sites avec remise,
  \item \alert{999 iterations},
  \item construction d'intervalle de confiance bootstrap \alert{90\%}.
\end{itemize}

\textbf{Notation:} BCI 90\% (Bootstrap Confidence Interval).
\end{frame}

\begin{frame}{Bootstrap: explication intuitive}
\textbf{Idee simple:}
\begin{itemize}
  \item on veut savoir si le signal regional change quand la composition des lacs change,
  \item le bootstrap simule cette variabilite en reconstruisant le regional plusieurs fois,
  \item chaque reconstruction tire \textit{6 lacs avec remise} parmi les 6 lacs disponibles.
\end{itemize}

\textbf{Avec remise, cela veut dire:}
\begin{itemize}
  \item un lac peut apparaitre plusieurs fois dans un tirage,
  \item un autre lac peut ne pas apparaitre dans ce tirage,
  \item cela imite l'incertitude d'echantillonnage inter-sites.
\end{itemize}

\textbf{Resultat:} on obtient non pas une seule courbe, mais une distribution de courbes possibles.
\end{frame}

\begin{frame}[fragile]{Bootstrap: pseudo-code}
\footnotesize
\begin{verbatim}
for b in 1..999:
    sample_lakes <- resample_with_replacement(6 lakes)
    compute RegBB_b(t), RegFF_b(t), RegFS_b(t)
    compute RegFRI_b(t) = 1 / RegFF_b(t)
end

for each time t:
    CI90_low(t)  = quantile(metric_b(t), 0.05)
    CI90_high(t) = quantile(metric_b(t), 0.95)
\end{verbatim}

\normalsize
\textbf{Regle de comparaison de periodes:}
\begin{itemize}
  \item difference consideree significative si les BCI 90\% ne se recouvrent pas.
\end{itemize}
\end{frame}

\begin{frame}{Comment lire un BCI 90\%?}
\textbf{A un temps \(t\), pour une metrique donnee:}
\begin{itemize}
  \item borne basse = quantile 5\%,
  \item borne haute = quantile 95\%,
  \item intervalle \([q_{0.05}, q_{0.95}]\) = plage plausible de la metrique.
\end{itemize}

\textbf{Interpretation pratique:}
\begin{itemize}
  \item intervalle large \(\Rightarrow\) forte incertitude inter-sites,
  \item intervalle etroit \(\Rightarrow\) bonne stabilite du signal regional.
\end{itemize}

\textbf{Comparaison de periodes:}
\begin{itemize}
  \item si les BCI 90\% se recouvrent fortement, difference prudente,
  \item s'ils ne se recouvrent pas, difference plus robuste.
\end{itemize}
\end{frame}

\begin{frame}{Que mettre dans la section lissage + bootstrap?}
\textbf{Lissage (obligatoire a rapporter):}
\begin{itemize}
  \item variable lissee (FF),
  \item methode (noyau),
  \item bande passante (500 ans),
  \item grille temporelle,
  \item traitement des bords.
\end{itemize}

\textbf{Bootstrap (obligatoire a rapporter):}
\begin{itemize}
  \item unite de resampling (site/lac),
  \item nombre d'iterations (999),
  \item niveau de confiance (90\%),
  \item quantiles utilises (5\%--95\%),
  \item critere de significativite (recouvrement des IC).
\end{itemize}
\end{frame}

\section{6. Parametrage pratique et controles}

\begin{frame}{Parametres cles a figer avant analyse}
\textbf{Bloc Chronologie:}
\begin{itemize}
  \item choix des dates retenues/exclues,
  \item prior/model age-profondeur,
  \item resolution temporelle d'interpolation.
\end{itemize}

\textbf{Bloc CharAnalysis:}
\begin{itemize}
  \item largeur de fenetre background,
  \item regle de seuil de pic,
  \item criteres de filtrage des faux pics.
\end{itemize}

\textbf{Bloc FF/FRI:}
\begin{itemize}
  \item noyau utilise,
  \item bande passante 500 ans,
  \item pas temporel des sorties.
\end{itemize}
\end{frame}

\begin{frame}{Controles qualite recommandes}
\begin{enumerate}
  \item \textbf{QC laboratoire}: replicats de comptage/aire de charbon.
  \item \textbf{QC chronologie}: verifier inversions/ruptures improbables age-profondeur.
  \item \textbf{QC pics}: test sensibilite aux parametres CharAnalysis.
  \item \textbf{QC regional}: leave-one-lake-out pour voir la dependance a un site.
  \item \textbf{QC incertitude}: stabilite des IC en augmentant le nombre d'iterations.
\end{enumerate}
\end{frame}

\begin{frame}{Erreurs frequentes a eviter}
\begin{itemize}
  \item Confondre \textit{flux de charbon} et \textit{nombre d'evenements}.
  \item Appliquer un lissage trop agressif et perdre les transitions rapides.
  \item Comparer des periodes sans verifier la taille de l'incertitude.
  \item Interpretrer RegFS comme une mesure directe de severite absolue.
  \item Oublier l'effet de la chronologie sur toutes les metriques derivees.
\end{itemize}
\end{frame}

\section{7. Reproductibilite}

\begin{frame}{Traite de reproductibilite (a inclure dans le rapport)}
\textbf{Checklist minimale:}
\begin{itemize}
  \item version des donnees (CSV/ZIP),
  \item version des outils (CharAnalysis, R/Python packages),
  \item scripts d'execution dans l'ordre exact,
  \item parametres numeriques fixes,
  \item sortie des tables intermediaires (CHAR, pics, FF, Reg*),
  \item archive des figures finales et metadonnees.
\end{itemize}
\end{frame}

\begin{frame}{Schema de flux complet}
\small
\textbf{Entrees}\
carottes sedimentaires + datations \(^{210}\)Pb/\(^{14}\)C

\vspace{0.15cm}
\textbf{Traitements}\
extraction >160 \(\mu\)m \(\rightarrow\) aire charbon \(\rightarrow\) CHAR \(\rightarrow\) CharAnalysis (pics) \\
\(\rightarrow\) FF (noyau 500 ans) \(\rightarrow\) agregation 6 lacs \(\rightarrow\) RegBB, RegFF, RegFS, RegFRI \\
\(\rightarrow\) bootstrap 999 (BCI 90\%)

\vspace{0.15cm}
\textbf{Sorties}\
courbes regionales + IC 90\% + comparaisons par periode
\end{frame}

\begin{frame}{References methodologiques a citer}
\footnotesize
\begin{itemize}
  \item Girardin et al. (2024), Communications Earth \& Environment, Methods section.
  \item Higuera et al. (2010), Peak detection in sediment-charcoal records.
  \item Brossier et al. (2014), calibration des pics avec tree-rings.
  \item Ali et al. (2009), area vs number vs volume in charcoal records.
  \item Higuera (2009), CharAnalysis v1.1.
\end{itemize}

\vspace{0.2cm}
\normalsize
\textbf{Message final:} la robustesse de l'interpretation depend autant du protocole CharAnalysis + bootstrap que de la qualite des chronologies age-profondeur.
\end{frame}

\section{8. Annexe - Notations scientifiques}

\begin{frame}{Pourquoi ces notations?}
\begin{itemize}
  \item Les articles paleo-environnementaux utilisent des notations compactes.
  \item Elles melangent isotopes, unites, statistiques et transformations.
  \item Cette annexe donne une traduction simple: \textit{notation -> sens -> exemple}.
\end{itemize}

\vspace{0.2cm}
\textbf{Objectif:} que chaque symbole vu dans la methode soit interpretable sans ambiguite.
\end{frame}

\begin{frame}{Isotopes et datations: \(^{210}\)Pb, \(^{14}\)C, AMS}
\textbf{\(^{210}\)Pb (plomb-210):}
\begin{itemize}
  \item isotope radioactif utilise pour dater les sediments recents (dernieres centuries),
  \item utile pour caler le sommet des carottes.
\end{itemize}

\textbf{\(^{14}\)C (carbone-14):}
\begin{itemize}
  \item isotope radioactif utilise pour dater les materiaux plus anciens,
  \item applique a des macro-restes (bois, feuilles, graines, etc.).
\end{itemize}

\textbf{AMS (Accelerator Mass Spectrometry):}
\begin{itemize}
  \item methode de mesure tres sensible du \(^{14}\)C,
  \item permet de dater de petites quantites de matiere.
\end{itemize}
\end{frame}

\begin{frame}{Notations de temps: BP, cal. yr BP, CE 1950}
\textbf{BP = Before Present:}
\begin{itemize}
  \item "Present" est fixe conventionnellement a l'annee \alert{1950 CE}.
\end{itemize}

\textbf{cal. yr BP = calibrated years before present:}
\begin{itemize}
  \item age calibre en annees avant 1950,
  \item "calibre" signifie corrige via courbe de calibration (ex. IntCal20).
\end{itemize}

\textbf{Exemples:}
\begin{itemize}
  \item 0 cal. yr BP = 1950 CE,
  \item 1000 cal. yr BP \(\approx\) 950 CE,
  \item 12000 cal. yr BP \(\approx\) 10050 BCE.
\end{itemize}
\end{frame}

\begin{frame}{IntCal20: definition courte}
\textbf{IntCal20 = courbe de calibration radiocarbone (version 2020).}

\vspace{0.2cm}
\textbf{A quoi elle sert?}
\begin{itemize}
  \item Une date \(^{14}\)C brute n'est pas directement une date calendrier.
  \item Le \(^{14}\)C atmospherique a varie dans le temps.
  \item IntCal20 corrige cette variabilite pour passer de \(^{14}\)C age BP a cal. yr BP.
\end{itemize}

\textbf{Message cle:}
\begin{itemize}
  \item sans calibration, les ages peuvent etre biaises,
  \item avec calibration, on obtient une plage d'ages calendrier probable.
\end{itemize}
\end{frame}

\begin{frame}{IntCal20 en 4 etapes (pratique)}
\begin{enumerate}
  \item Le labo fournit une date radiocarbone: \(x \pm \sigma\) \(^{14}\)C BP.
  \item On utilise IntCal20 pour convertir cet age en age calendrier.
  \item On obtient une distribution de probabilite (pas une valeur unique).
  \item On resume cette distribution par des intervalles (ex. 68\%, 95\%) en cal. yr BP.
\end{enumerate}

\vspace{0.2cm}
\textbf{Formulation simple:}
\[
\text{age calendrier} = f\big(\text{age }^{14}\text{C},\, \text{IntCal20}\big)
\]
\end{frame}

\begin{frame}{Exemple numerique (pedagogique)}
\textbf{Exemple fictif (ordre de grandeur):}
\begin{itemize}
  \item date labo: \(5000 \pm 40\) \(^{14}\)C BP,
  \item apres calibration IntCal20: distribution centree autour de ~5.7--5.9 ka cal BP.
\end{itemize}

\textbf{Important:}
\begin{itemize}
  \item ce n'est pas une simple addition/soustraction,
  \item la conversion depend de la forme de la courbe IntCal20,
  \item un meme age \(^{14}\)C peut donner une distribution asymetrique ou multimodale.
\end{itemize}

\textbf{Dans ce projet:}
\begin{itemize}
  \item les dates calibrees (cal. yr BP) alimentent le modele age-profondeur,
  \item puis tout le reste (CHAR, pics, FF, FRI, RegFS) en depend.
\end{itemize}
\end{frame}

\begin{frame}{Unites charbon: \(\mu\)m, mm\(^2\) cm\(^{-3}\), CHAR}
\textbf{\(\mu\)m (micrometre):}
\begin{itemize}
  \item 1 \(\mu\)m = 10\(^{-6}\) m,
  \item seuil >160 \(\mu\)m = fragments macroscopiques retenus.
\end{itemize}

\textbf{mm\(^2\) cm\(^{-3}\):}
\begin{itemize}
  \item aire de charbon mesuree par volume de sediment,
  \item variable "brute" avant conversion temporelle.
\end{itemize}

\textbf{CHAR en mm\(^2\) cm\(^{-2}\) an\(^{-1}\):}
\begin{itemize}
  \item flux de charbon par surface et par an,
  \item comparable dans le temps et entre lacs.
\end{itemize}

\textbf{Exemple:} CHAR = 0.02 mm\(^2\) cm\(^{-2}\) an\(^{-1}\)
\begin{itemize}
  \item signifie qu'en moyenne, 0.02 mm\(^2\) de charbon se depose chaque annee sur 1 cm\(^2\).
\end{itemize}
\end{frame}

\begin{frame}{Notations de metriques feu: FF, FRI, RegBB, RegFS}
\textbf{FF (fire frequency):}
\begin{itemize}
  \item frequence des evenements feu (feux/an),
  \item estimee apres detection de pics + lissage noyau.
\end{itemize}

\textbf{FRI (fire return interval):}
\[
FRI = \frac{1}{FF}
\]

\textbf{Exemple:}
\begin{itemize}
  \item si FF = 0.005 feux/an, alors FRI = 200 ans.
\end{itemize}

\textbf{RegBB, RegFF, RegFS, RegFRI:}
\begin{itemize}
  \item prefixe "Reg" = moyenne/indice regional (agregation multi-sites),
  \item RegFS = RegBB / RegFF (proxy taille-severite moyenne par evenement).
\end{itemize}
\end{frame}

\begin{frame}{Notations statistiques: CI, BCI, 90\%, 999 iterations}
\textbf{CI (Confidence Interval):}
\begin{itemize}
  \item intervalle d'incertitude autour d'une estimation.
\end{itemize}

\textbf{BCI (Bootstrap Confidence Interval):}
\begin{itemize}
  \item CI calcule par re-echantillonnage bootstrap.
\end{itemize}

\textbf{IC 90\%:}
\begin{itemize}
  \item contient 90\% des valeurs bootstrap (souvent quantiles 5\% et 95\%).
\end{itemize}

\textbf{999 iterations:}
\begin{itemize}
  \item 999 re-echantillonnages des lacs pour estimer la variabilite regionnale.
\end{itemize}

\textbf{Exemple rapide:}
\begin{itemize}
  \item RegFRI median = 210 ans,
  \item BCI 90\% = [180 ; 245] ans.
\end{itemize}
\end{frame}

\begin{frame}{Mini glossaire (resume en 60 secondes)}
\small
\begin{tabular}{ll}
\toprule
\textbf{Notation} & \textbf{Definition courte} \\
\midrule
\(^{210}\)Pb & datation sediments recents \\
\(^{14}\)C & datation plus ancienne \\
cal. yr BP & annees calibrees avant 1950 \\
CHAR & flux de charbon (surface \& temps) \\
FF & frequence de feux \\
FRI & intervalle moyen entre feux (1/FF) \\
RegBB & biomasse brulee regionale \\
RegFS & indice taille/severite regional \\
CI / BCI 90\% & incertitude (intervalle 90\%) \\
\bottomrule
\end{tabular}
\end{frame}

\begin{frame}
\centering
\Huge Merci
\end{frame}

\end{document}
